% !TeX spellcheck = es_ES

% Comandos para crear la portada y la primera página del PFG
%% Título del PFG
\newcommand{\PFGtitulo}{Engine para juegos de tipo Game \& Watch.\\ 
	{\Large Diario de iteraciones del proyecto}}
%% Título corto, para mostrar en el encabezado de las páginas pares. Por defecto se pone el título, pero, si es muy largo, puede acortarse aquí
\newcommand{\PFGtitulocorto}{Iteraciones}
%% Autor/a y directores/Fs del PFG
\newcommand{\PFGautor}{Jose Antonio Vaqué Urbaneja}
\newcommand{\PFGdirector}{Alberto Pérez}
%% Descomentar y rellenar la siguiente línea si hay codirector/a en el PFG
\newcommand{\PFGcodirector}{Fernando López Ostenero}
%% Grado. A elegir entre:
%% Ingeniería Informática
%% Ingeniería en Tecnologías de la Información
\newcommand{\PFGgrado}{Ingeniería Informática}
%% Modalidad. A elegir entre:
%% genérica
%% específica
\newcommand{\PFGmodalidad}{específica}
%% Curso académico en formato 20xx/20xx+1. Por ejemplo 2022/2023
\newcommand{\PFGcursoacademico}{2026/2027}
%% Convocatoria. A elegir entre:
%% junio
%% septiembre
%% diciembre
\newcommand{\PFGconvocatoria}{junio}
%% Fecha de lectura (si no se sabe el día concreto, poner mes y año}
\newcommand{\PFGfechadefensa}{junio 2027}


\documentclass[12pt,twoside]{book}
\usepackage{graphicx}
\usepackage[utf8]{inputenc}
\usepackage[papersize={210mm,297mm},tmargin=25mm,bmargin=25mm,lmargin=20mm,rmargin=20mm]{geometry}

% No indentar los párrafos y separarlos 1cm
\setlength\parindent{0cm}
\setlength\parskip{1em}

% Para tener los elementos de texto en español
\usepackage[spanish]{babel}
% Para poder resaltar un texto en amarillo con \hl{}
\usepackage{soul}
% Definición de colores
\usepackage{xcolor}
\definecolor{links}{rgb}{0.35, 0.4, 0.35}
\definecolor{urls}{rgb}{0,0,1}
\definecolor{cites}{rgb}{0.8,0,0}
\definecolor{gray20}{gray}{0.20}

% Hiperreferencias
\usepackage[hyphens]{url}
\usepackage{hyperref}
\hypersetup{
	colorlinks = true,
	linkcolor = links,
	urlcolor = urls,
	citecolor = cites
}

%%%%%%%%%%%%%%%%%%%%%%%%%%%%%%%%%%%%%%%%%%%%%%%%%%%%%%%%%%%%%%%%%%%%%%%%%%%%%%%%%%%%%%%%%%%%%%%
%%%%%%%%%%%%%%%%% PAQUETES AÑADIDOS %%%%%%%%%%%%%%%%%%%%%%%%%%%%%%%%%%%%%%%%%%%%%%%%%%%%%%%%%%%
%%%%%%%%%%%%%%%%%%%%%%%%%%%%%%%%%%%%%%%%%%%%%%%%%%%%%%%%%%%%%%%%%%%%%%%%%%%%%%%%%%%%%%%%%%%%%%%
\usepackage{ltablex} % Combina las características de tabularx y longtable
\usepackage{booktabs} % Para líneas de tabla más limpias y profesional
\usepackage{booktabs,xltabular}
\usepackage{float}
\usepackage{amsmath}
\usepackage{textcomp}
\usepackage{makecell}
\usepackage{dirtree}
\usepackage{ragged2e}
\usepackage[normalem]{ulem}
\usepackage{needspace}
\usepackage[spanish, onelanguage, ruled, vlined, linesnumbered]{algorithm2e}
\IncMargin{1.5em} % Empuja todo el contenido (números incluidos) hacia la derecha
% El primer valor es el espacio antes de la línea vertical, el segundo es la sangría del código
\SetInd{1.5em}{1.5em}
\SetNlSkip{0.8em} % <-- Separa los números de línea del código
\SetKwProg{Fn}{Función}{:}{fin}

% Cambiamos títulos de los capítulos
\addto\captionsspanish{%
	\renewcommand{\chaptername}{FASE}%
	\renewcommand{\appendixname}{ANEXO}%
}
%%%%%%%%%%%%%%%%%%%%%%%%%%%%%%%%%%%%%%%%%%%%%%%%%%%%%%%%%%%%%%%%%%%%%%%%%%%%%%%%%%%%%%%%%%%%%%%



% Formato de los listados
\usepackage{listings}
\lstloadlanguages{C,C++,Haskell,HTML,Java,PHP,Python}

\definecolor{codegreen}{rgb}{0,0.6,0}
\definecolor{codegray}{rgb}{0.5,0.5,0.5}
\definecolor{codepurple}{rgb}{0.58,0,0.82}
\definecolor{backcolour}{rgb}{0.95,0.95,0.92}

\lstdefinestyle{mystyle}{
    commentstyle=\color{codegreen},
    keywordstyle=\color{blue},
    numberstyle=\tiny\color{codegray},
    stringstyle=\color{codepurple},
    backgroundcolor=\color{backcolour}, % choose the background color; you must add \usepackage{color} or \usepackage{xcolor}; should come as last argument
	basicstyle=\footnotesize\ttfamily\color{black},        % the size of the fonts that are used for the code
	belowcaptionskip=1\baselineskip,
	breakatwhitespace=false,         % sets if automatic breaks should only happen at whitespace
	breaklines=true,                 % sets automatic line breaking
	captionpos=b,                    % sets the caption-position to bottom
	deletekeywords={...},            % if you want to delete keywords from the given language
	escapeinside={\%*}{*)},          % if you want to add LaTeX within your code
	extendedchars=true,              % lets you use non-ASCII characters; for 8-bits encodings only, does not work with UTF-8
	firstnumber=1,                   % start line enumeration with line 1000
	frame=single,
	%identifierstyle=\color{blue!50!black},
	keepspaces=true,                 % keeps spaces in text, useful for keeping indentation of code (possibly needs columns=flexible)
	%language=C,                     % the language of the code
	%morekeywords={*,...},           % if you want to add more keywords to the set
	numbers=left,                    % where to put the line-numbers; possible values are (none, left, right)
	numbersep=10pt,                  % how far the line-numbers are from the code
	numberstyle=\tiny\color{gray},   % the style that is used for the line-numbers
	rulecolor=\color{black},         % if not set, the frame-color may be changed on line-breaks within not-black text (e.g. comments (green here))
	morecomment=[n]{(*}{*)},
	morecomment=[n][\color{magenta}]{\%I}{\ },
	morecomment=[n][\color{magenta}]{\%Q}{\ },
	showspaces=false,                % show spaces everywhere adding particular underscores; it overrides 'showstringspaces'
	showstringspaces=false,          % underline spaces within strings only
	showtabs=false,                  % show tabs within strings adding particular underscores
	stepnumber=1,                    % the step between two line-numbers. If it's 1, each line will be numbered
	tabsize=2,	                   % sets default tabsize to 2 spaces
	title=\lstname,
	xleftmargin=1.5cm,
}
\lstset{style=mystyle}

% Definición de colores
\definecolor{jsonpurple}{rgb}{0.5,0,0.5}
\definecolor{jsoncyan}{rgb}{0,0.6,0.6}
\definecolor{jsonorange}{rgb}{0.8,0.4,0}
\definecolor{jsongray}{rgb}{0.4,0.4,0.4}

% Definición CORRECTA del lenguaje JSON para listings
\lstdefinelanguage{json}{
	stringstyle=\color{jsonpurple},
	morestring=[b]",
	morestring=[d]',
	literate=
	*{0}{{{\color{jsonorange}0}}}{1}
	{1}{{{\color{jsonorange}1}}}{1}
	{2}{{{\color{jsonorange}2}}}{1}
	{3}{{{\color{jsonorange}3}}}{1}
	{4}{{{\color{jsonorange}4}}}{1}
	{5}{{{\color{jsonorange}5}}}{1}
	{6}{{{\color{jsonorange}6}}}{1}
	{7}{{{\color{jsonorange}7}}}{1}
	{8}{{{\color{jsonorange}8}}}{1}
	{9}{{{\color{jsonorange}9}}}{1}
	{:}{{{\color{jsoncyan}:}}}{1}
	{,}{{{\color{jsoncyan},}}}{1}
	{\{}{{{\color{jsongray}{\{}}}}{1}
	{\}}{{{\color{jsongray}{\}}}}}{1}
	{[}{{{\color{jsongray}{[}}}}{1}
	{]}{{{\color{jsongray}{]}}}}{1},
}

% Añadir ficheros externos
\usepackage{subfiles}
% Soporte para anexos
\usepackage{appendix}
% Paquetes para figuras
\usepackage{caption}
\usepackage{subcaption}
% Para las tablas
\usepackage{tabularx}
\usepackage{multicol}
\usepackage{multirow}
%Encabezados bonitos y acordes al reglamento
\usepackage{fancyhdr}
\fancyhf{}
\lhead[\thepage]{\leftmark}
\rhead[\PFGtitulocorto]{\thepage}
\setlength{\headheight}{16pt}
\pagestyle{empty}

% Paquete para poner la bibliografia en la table of contents
\usepackage[nottoc,notlot,notlof]{tocbibind}
% Cambiamos el espacio en blanco enorme que por defecto hay en las páginas inicio de capítulo
\usepackage{titlesec}
\titleformat{\chapter}[display]
{\normalfont\huge\bfseries}{\chaptertitlename\ \thechapter}{20pt}{\Huge}
\titlespacing*{\chapter}{0pt}{-50pt}{40pt}

% Para crear subsubsubsection
\setcounter{secnumdepth}{4}
\titleformat{\paragraph}
{\normalfont\normalsize\bfseries}{\theparagraph}{1em}{}
\titlespacing*{\paragraph}
{0pt}{3.25ex plus 1ex minus .2ex}{1.5ex plus .2ex}

% Ruta en la que pondremos las imágenes para el documento
\graphicspath{ {imagenes/} }
% Redifinimos el nombre que asigna babel spanish a las tablas (cuadro) a tabla
\renewcommand{\spanishtablename}{Tabla}
\renewcommand{\spanishlisttablename}{Índice de tablas}
\renewcommand{\spanishcontentsname}{Índice}
% Redefinimos el nombre de apéndice a anexo
\renewcommand{\appendixname}{Anexos}
\renewcommand{\appendixtocname}{Anexos}
\renewcommand{\appendixpagename}{Anexos}

% Definimos el comando de página en blanco
\newcommand\paginablanco{%
    \null
    \thispagestyle{empty}%
    \newpage}

% Definimos el comando de página en blanco sin avanzar numeracion
\newcommand\paginablancosin{%
    \paginablanco{}
    \addtocounter{page}{-1}}

% Mayor anchura en las tablas
\renewcommand{\arraystretch}{1.5}
% Mayor ancuhura en las fracciones
\newcommand\ddfrac[2]{\frac{\displaystyle #1}{\displaystyle #2}}
% Comando escribir PFG
\newcommand{\pfg}{Proyecto Fin de Grado }


% Comenzamos el documento
\begin{document}
\pagenumbering{gobble} % para no incluir números en la portada y primera página

% Hacemos la portada
\begin{titlepage}

\begin{figure}[!htb]
    \centering
    \includegraphics{rosadelosvientos.png}
\end{figure}

\begin{center}
    {\scshape\Large Universidad Nacional de Educación a Distancia \par}
    {\scshape\Large Escuela Técnica Superior de Ingeniería Informática \par}
    \vfill
    \large{Proyecto de fin de Grado en \PFGgrado}\\
    \vfill
    {\bfseries\Huge \textcolor{black}{\PFGtitulo} \par}
    \vfill
\end{center}

\vspace*{3cm}
{\huge \textcolor{black}{\PFGautor} \par}
{\huge \textcolor{black}{Dirigido por: \PFGdirector} \par}
\ifdefined\PFGcodirector
  \huge \textcolor{black}{Codirigido por: \PFGcodirector} \par
\fi
{\huge \textcolor{black}{Curso \PFGcursoacademico, convocatoria \PFGconvocatoria} \par}

\chapter*{}

\thispagestyle{empty}

\begin{figure}[!htb]
    \centering
    \includegraphics{rosadelosvientos.png}
\end{figure}

\begin{center}

    {\bfseries\Huge \textcolor{black}{\PFGtitulo} \par}
    \vfill
    {\bfseries\Large \textcolor{black}{Proyecto de fin de Grado en \PFGgrado} \par}
    {\bfseries\Large \textcolor{black}{de modalidad \PFGmodalidad} \par}
\end{center}

\vspace*{5cm}
{\Large Realizado por: \PFGautor \par}
{\Large Dirigido por: \PFGdirector \par}
\ifdefined\PFGcodirector
  {\Large Codirigido por: \PFGcodirector \par}
\fi%

%%{\Large Codirigido por: \PFGcodirector \par}
\vfill
{\Large Fecha de lectura y defensa: \PFGfechadefensa \par}

\end{titlepage}

\pagenumbering{Roman} % para comenzar la numeracion de paginas en numeros romanos

\chapter*{Resumen} % El * indica que no queremos que añada la palabra "Capitulo"

Este documento es el diario de las iteraciones efectuadas para el desarrollo del proyecto, consistente en implementar un motor que permita crear juegos de tipo \textbf{Game \& Watch}. Em el documento principal está el resumen de cada Fase, este es un documento anexo con el detalle que he ido anotando en cada una de las iteraciones de cada fase.

\chapter*{Abstract} % El * indica que no queremos que añada la palabra "Capitulo"

This document is the journal of the iterations carried out for the development of the project, which consists of implementing an engine to create Game \& Watch-type games. The main document contains the summary of each Phase; this is an attached document with the details that I have been noting down during each of the iterations of each phase..

\tableofcontents %% Índice general

%%%%%%%%%%%%%%%%%%%%%%%%%%%%%%%%%%%%%%%%%%%%%%%%%%%%%%%%%%%%%%%%%%%%%%%%%%%%%%%%%%%%%%%%%%%%%%%
%%%%%%%%%%%%%%%%% AÑADIDO %%%%%%%%%%%%%%%%%%%%%%%%%%%%%%%%%%%%%%%%%%%%%%%%%%%%%%%%%%%%%%%%%%%%%
%%%%%%%%%%%%%%%%%%%%%%%%%%%%%%%%%%%%%%%%%%%%%%%%%%%%%%%%%%%%%%%%%%%%%%%%%%%%%%%%%%%%%%%%%%%%%%%
% Redifinimos el nombre que asigna babel spanish a las tablas (cuadro) a tabla
\renewcommand{\lstlistingname}{Código}
\renewcommand{\lstlistlistingname}{Índice de códigos}
%%%%%%%%%%%%%%%%%%%%%%%%%%%%%%%%%%%%%%%%%%%%%%%%%%%%%%%%%%%%%%%%%%%%%%%%%%%%%%%%%%%%%%%%%%%%%%%

\lstlistoflistings % <-- Añade aquí el Índice de Códigos Fuentes
\addcontentsline{toc}{section}{Índice de listados} % Queremos que aparezca en el índice


\newpage

\paginablanco{}

%% Comenzamos con la numeración arábiga y el estilo de cabeceras de páginas definido, porque aquí empieza el contenido
\pagenumbering{arabic}
\pagestyle{fancy}
%% Importación de capítulos
\subfile{capitulos/F01}
\subfile{capitulos/F02}
\subfile{capitulos/F03}
\subfile{capitulos/F04}
\subfile{capitulos/F05}
\subfile{capitulos/F06}
\subfile{capitulos/F07}
\appendix
\subfile{capitulos/A01 Funciones}
\subfile{capitulos/A02 Formatos}

%...

\end{document}